\documentclass[11pt]{article}

\usepackage[margin=1in]{geometry}
\usepackage{microtype}
\usepackage{amsmath}
\usepackage{graphicx}
\usepackage{booktabs}
\usepackage[hidelinks]{hyperref}

\graphicspath{{figures/}}

\newcommand{\natex}{\texttt{natex}}

\title{A Genuine Bend, No Causal Certificate: Difference-in-Kinks Tests of
Hyperscaler Capital Expenditure at ChatGPT}
\author{Hauke Hillebrandt\thanks{University College London. Email:
\texttt{ucjthhi@ucl.ac.uk}. This paper and the underlying \natex{} software were
prepared with substantial assistance from Anthropic's Claude models; the author
reviewed the analyses and text and is responsible for all remaining errors.}}
\date{July 2026 \\ \small analysis pass of record: 2026-07, \natex{} v0.2.0, seed 0}

\begin{document}

\maketitle

\begin{abstract}
\noindent Quarterly capital expenditure of the big-4 hyperscalers (Microsoft,
Alphabet, Amazon, Meta) accelerated sharply relative to a non-AI
capital-intensive comparison aggregate after the release of ChatGPT on 30
November 2022. A group difference-in-kinks (DiK) estimate at the release date
is $+0.1396$ dex/yr (HC1 se $0.0302$, nominal $p=3.9\times10^{-6}$), positive
in all 18 specification cells ($+0.049$ to $+0.195$), localized over
2023.0--2023.5 exactly as a lagged capex response to a late-2022 shock would
be, and present per firm for Microsoft ($+0.171$), Alphabet ($+0.206$), and
Amazon ($+0.137$) but null for Meta ($+0.031$, $p=0.55$). Attribution to
ChatGPT, however, fails the falsification battery: a clean pre-period
placebo-cutoff grid rejects at nominal 5\% in 7 of 15 non-event dates
(empirical size $0.47$), giving placebo-calibrated $p=0.125$ ($|z|$) and
$0.250$ ($|\tau|$); the largest placebo $|z|$ ($5.11$, at the non-event date
2019.75) exceeds the ChatGPT estimate's matched-bandwidth $|z|$ ($4.35$); a
full-sample kink at the non-event date 2021.5 ($-0.314$, $|z|=6.9$) is larger
than the main effect; and a cutoff at the BIS semiconductor export controls
eight weeks earlier is statistically indistinguishable ($+0.125$, $z=4.0$).
Residual lag-1 autocorrelation reaches $0.66$, so the nominal
heteroskedasticity-robust $p$-values are artifacts of applying
serial-independence inference to smooth aggregates. The slope divergence is a
genuine, robust data feature; no causal attribution to ChatGPT survives
triage.
\end{abstract}

\section{Introduction}\label{sec:intro}

OpenAI released ChatGPT on 30 November 2022. Within two years the capital
expenditure of the four large US ``hyperscalers'' --- Microsoft, Alphabet,
Amazon, and Meta --- had become a macroeconomic quantity in its own right,
central to debates about how much aggregate investment the current AI wave can
move \cite{acemoglu2025}. Financial markets repriced AI exposure almost
immediately: \cite{eisfeldt2023} construct firm-level workforce exposures to
generative AI and show that an exposure-sorted long-short portfolio earned
roughly 5\% in the two weeks after the release. A natural real-economy
question follows: did the release bend the hyperscalers' \emph{investment}
trend --- the growth rate of measured capital expenditure --- or was the
buildout already accelerating?

Any calendar-time answer must contend with a crowded event window. Eight weeks
before ChatGPT, on 7 October 2022, the Bureau of Industry and Security
announced sweeping export controls on advanced AI chips and semiconductor
manufacturing equipment, a policy landmark in its own right \cite{allen2022};
any quarterly series treats the two events as nearly the same date.

The estimand --- a change in trend slope for one group relative to another at a
known date --- is the difference-in-kinks (DiK) design formalized by
\cite{bockerman2025}, building on the regression kink design of
\cite{card2015}. Kink estimates on time-series running variables are known to
produce spuriously significant slope changes, which is why placebo
distributions at non-event dates are the appropriate yardstick
\cite{ganong2018}; and conventional robust standard errors are badly oversized
on serially correlated aggregates \cite{bertrand2004}. This paper applies the
\natex{} toolkit's DiK estimator \cite{natex} to a two-group quarterly capex
panel built from SEC EDGAR filings, with ChatGPT as the candidate cutoff and a
mandatory falsification battery --- which converts a nominally overwhelming
event-study result into an explicit non-attribution.

\section{Data}\label{sec:data}

Quarterly capital expenditure is reconstructed from SEC EDGAR XBRL
\texttt{companyconcept} facts, converting fiscal year-to-date values to clean
quarters by same-start differencing with remainder assignment across tag
transitions. The tag is
\texttt{PaymentsToAcquirePropertyPlantAndEquipment}; Amazon additionally
reports under \texttt{PaymentsToAcquireProductiveAssets}. The treated series is the big-4 aggregate over 46
quarters, 2014Q4--2026Q1, rising from \$6.8bn to \$129.8bn per quarter. The
rebuilt aggregate matches the frozen reference series of the run of record on
all 45 common quarters (maximum difference \$0.000000bn) and its processed
$\log_{10}$ counterpart on all 46 (maximum difference $5\times10^{-5}$).

The comparison series aggregates six large non-AI capital-intensive firms:
Walmart, ExxonMobil, AT\&T, Verizon, Union Pacific, and Home Depot, fetched
through the identical EDGAR pipeline. Firm-specific tag quirks (AT\&T,
Verizon, and Home Depot report quarterly capex under
\texttt{ProductiveAssets}-family tags; Walmart and Home Depot have
January-31 fiscal year ends mapped to the nearest calendar quarter) are
resolved against the live EDGAR records, and 2023 calendar-year sums match
published figures for all six firms (AT\&T \$17.85bn, Verizon \$18.77bn,
ExxonMobil \$21.92bn, Walmart \$20.61bn, Union Pacific \$3.61bn, Home Depot
\$3.23bn).

The outcome is $\log_{10}$ of group capex in \$bn per quarter, so slopes are
in dex/yr (one dex is a factor of ten). The running variable is the quarter
midpoint in fractional years; the ChatGPT cutoff is $t_0=2022.913$
(2022-11-30). A donut of $0.13$ years ($\approx$ one quarter) removes 2022Q4,
the quarter that straddles the release.

\section{Design and Methods}\label{sec:methods}

\paragraph{Group difference-in-kinks.} Following \cite{bockerman2025}, the DiK
estimand is the change in the outcome's slope in calendar time at the cutoff
for the treated group \emph{minus} the same slope change for the comparison
group, estimated by local linear weighted regression on each side of the
cutoff within each group (\natex{} v0.2.0 \cite{natex}; the group indicator
enters the CLI's DiK period dimension with a unit policy-kink change, so
$\hat\tau = (\text{kink})_{\text{big-4}} - (\text{kink})_{\text{control}}$ in
dex/yr). The primary specification uses bandwidth 3.0 years, triangular
kernel, the one-quarter donut, and HC1 standard errors; $n=46$ quarter-cells
(11 pre-cutoff and 12 post-cutoff quarters per group).

\paragraph{Honest-inference caveats, stated as run.} The estimator's own
caveats apply verbatim: the cutoff and bandwidth are user-specified, so
bandwidth, donut, and placebo-cutoff sensitivity must be reported; the
reported interval is conventional local-polynomial inference and may retain
smoothing bias; and causal interpretation requires parallel changes in
non-policy slope kinks and a time-stable marginal response at the cutoff.
HC1/CR1 inference additionally assumes serial independence, which these
aggregates violate: residual lag-1 autocorrelation in the primary fit is
$0.657$ on the control pre-cutoff side and $0.367$ on the big-4 pre-cutoff
side, precisely the setting in which nominal robust standard errors are known
to be understated \cite{bertrand2004}.

\paragraph{Falsification battery.} (i) an 18-cell specification grid
(bandwidth $\in\{2,3,4\}$ years $\times$ donut $\in\{0,0.13,0.25\}$
$\times$ \{triangular, uniform\}); (ii) a placebo cutoff at the BIS export
controls (2022-10-07, $t=2022.767$); (iii) a clean \emph{pre-period}
placebo-cutoff grid --- sample restricted to $t<2022.7$, 15 cutoffs
2017.0--2020.5, bandwidth 2.0, donut $0.13$ --- in the spirit of the
permutation logic of \cite{ganong2018}, yielding a placebo-calibrated
$p$-value for the main effect at the matched bandwidth; (iv) a cutoff
localization scan over 2021.5--2024.75; (v) per-firm staggered DiKs of each
of the ten firms against the control aggregate; and (vi) robustness checks:
CR1 clustering by year and by quarter, quarter-of-year dummies,
leave-one-out and no-retail control aggregates, a role-swap sanity check, and
a degree-2 polynomial.

\paragraph{Automated pipeline, stated as run.} The \natex{} survey pipeline
(\texttt{natex survey --seed 0}) was recorded for completeness: its automated
kink family pools both groups on the auto-picked linear capex outcome and
therefore cannot express the group-DiK aliasing; it returns null (Holm
$p=0.24$). The remaining families were skipped, needs-input, or failed as
expected on a 92-row two-unit panel. The manual DiK runs above are the
primary evidence.

\section{Results}\label{sec:results}

\begin{figure}[t]
\centering
\includegraphics[width=\textwidth]{fig1}
\caption{(a) Quarterly capital expenditure (\$bn, log scale) of the big-4
hyperscaler aggregate and the non-AI capital-intensive control aggregate,
2014Q4--2026Q1, with the BIS export-control (2022-10-07, dashed) and ChatGPT
(2022-11-30, solid) dates marked. (b) Group DiK estimates with 95\% HC1
confidence intervals at 15 pre-period non-event placebo cutoffs (gray;
pre-only sample, $t<2022.7$) and at the ChatGPT cutoff (vermilion), all at
the matched bandwidth 2.0 yr, triangular kernel, one-quarter donut; filled
markers are nominally significant at the 5\% level. Seven of 15 placebos
reject, and the largest placebo $|z|$ exceeds the ChatGPT $|z|$.}
\label{fig:fig1}
\end{figure}

\paragraph{The bend is real \ldots} The primary group DiK at ChatGPT
(Table~\ref{tab:dik}, first row) is $\hat\tau=+0.1396$ dex/yr (HC1 se
$0.0302$, 95\% CI $[0.0804, 0.1988]$, $z=4.62$, nominal
$p=3.85\times10^{-6}$). It decomposes into a big-4 kink of $+0.0878$ dex/yr
(slope $0.0999\to0.1876$; the aggregate's growth roughly doubles, to a pace
that multiplies capex by ten in about 5.3 years) and a control kink of
$-0.0518$ ($0.0811\to0.0293$). The estimate is positive in every cell of the
18-cell grid ($+0.0493$ to $+0.1951$, all nominal $p<0.05$; the uniform
kernel roughly halves $\hat\tau$ relative to the triangular at bandwidths
$\geq3$). The point estimate is essentially invariant to CR1 clustering by
year (se $0.0313$) or quarter (se $0.0278$), quarter-of-year dummies
($+0.1396$, se $0.0296$), leave-one-out control aggregates ($+0.1329$ to
$+0.1515$), and a no-retail control of ExxonMobil, AT\&T, Verizon, and Union
Pacific ($+0.1453$, se $0.0361$); the role-swap check returns exactly
$-0.1396$; the degree-2 estimate ($+0.3022$, se $0.0914$) shows curvature
sensitivity in magnitude but not sign. Per firm against the control
aggregate: Microsoft $+0.1710$ (se $0.0366$), Alphabet $+0.2061$ (se
$0.0349$), Amazon $+0.1375$ (se $0.0512$), Meta $+0.0308$ (se $0.0518$,
$p=0.55$ --- a true null: Meta's surge came in 2024, after its 2022--23
pullback); control firms' own kinks are all mildly negative ($-0.022$ to
$-0.090$). A localization scan (bandwidth 2.0, no donut) finds the DiK
crossing zero near $t=2022.3$, maximizing at $t=2023.25$ ($+0.2213$, se
$0.0440$), staying above $+0.19$ over 2023.0--2023.75, and decaying to
$\approx0$ by 2024.5 --- consistent with a lagged capex response to a
late-2022 shock.

\paragraph{\ldots\ but attribution to ChatGPT fails.} Three findings block a
causal reading. \emph{First}, the pre-period placebo battery fails
decisively: at 15 non-event cutoffs in 2017.0--2020.5, the identical
specification rejects at nominal 5\% seven times (empirical size $0.467$;
Figure~\ref{fig:fig1}b). Calibrated against this placebo distribution, the
ChatGPT estimate at the matched bandwidth ($+0.1778$, $z=4.35$) has
$p_{|z|}=(1+1)/(15+1)=0.125$ and $p_{|\tau|}=(3+1)/16=0.250$; the largest
placebo $|z|$ is $5.11$ (at 2019.75) and the largest placebo $|\hat\tau|$ is
$0.284$ (at 2019.5) --- both exceeding the ChatGPT values. \emph{Second}, the
full-sample kink at the non-event date 2021.5 is $-0.3135$ (se $0.0452$,
$|z|=6.9$), larger in magnitude and nominal significance than the main
effect: these smooth aggregates generate huge nominal kinks where nothing
happened. \emph{Third}, the cutoff is confounded: a DiK at the BIS
export-control date eight weeks earlier returns $+0.1253$ (se $0.0317$,
$z=3.95$), statistically indistinguishable from the ChatGPT estimate ---
quarterly data cannot separate the two events (or any other late-2022
candidate shock).

\begin{table}[t]
\centering
\small
\caption{Group difference-in-kinks estimates, $\log_{10}$ capex (dex/yr).}
\label{tab:dik}
\begin{tabular}{lrrrc}
\toprule
Specification & $\hat\tau$ & se & $|z|$ & sig.\ 5\% \\
\midrule
\multicolumn{5}{l}{\emph{Panel A: ChatGPT cutoff ($t_0=2022.913$)}}\\
triangular, bw 3, donut (primary)      & $+0.1396$ & $0.0302$ & $4.62$ & $\ast$ \\
triangular, bw 2, donut (placebo-matched) & $+0.1778$ & $0.0409$ & $4.35$ & $\ast$ \\
uniform, bw 4, donut (grid minimum)    & $+0.0493$ & $0.0210$ & $2.35$ & $\ast$ \\
primary, CR1 by year                   & $+0.1396$ & $0.0313$ & $4.46$ & $\ast$ \\
primary, quarter-of-year dummies       & $+0.1396$ & $0.0296$ & $4.72$ & $\ast$ \\
no-retail control (XOM, T, VZ, UNP)    & $+0.1453$ & $0.0361$ & $4.03$ & $\ast$ \\
degree 2, bw 3, donut                  & $+0.3022$ & $0.0914$ & $3.30$ & $\ast$ \\
\midrule
\multicolumn{5}{l}{\emph{Panel B: placebo cutoffs}}\\
BIS export controls, $t_0=2022.767$, bw 3 & $+0.1253$ & $0.0317$ & $3.95$ & $\ast$ \\
$t_0=2021.5$ (no event), full sample, bw 2 & $-0.3135$ & $0.0452$ & $6.93$ & $\ast$ \\
$t_0=2019.75$ (no event), pre-only, bw 2   & $+0.2627$ & $0.0514$ & $5.11$ & $\ast$ \\
$t_0=2019.5$ (no event), pre-only, bw 2    & $+0.2843$ & $0.0736$ & $3.86$ & $\ast$ \\
\bottomrule
\end{tabular}

\medskip
\parbox{0.92\textwidth}{\footnotesize Notes: HC1 standard errors against
standard normal critical values except where noted; donut $=0.13$ yr (one
quarter); bandwidths in years. Outcome: $\log_{10}$ group capex, \$bn per
quarter. The pre-only placebo battery (sample $t<2022.7$, 15 cutoffs
2017.0--2020.5, bw 2) rejects at nominal 5\% in 7/15 cases; the
placebo-calibrated $p$-value of the ChatGPT estimate at the matched bandwidth
is $0.125$ ($|z|$) and $0.250$ ($|\tau|$).}
\end{table}

\section{Caveats and Conclusion}\label{sec:conclusion}

Three caveats are decisive. First, the nominal inference is an artifact:
HC1/CR1 standard errors assume serial independence, but the outcome is a pair
of smooth, serially correlated aggregates (residual lag-1 autocorrelation up
to $0.66$), and the placebo battery measures the resulting size distortion
directly --- a 5\% test rejects 47\% of the time at non-event dates, the
failure mode documented by \cite{bertrand2004} and exactly what placebo
distributions exist to catch \cite{ganong2018}. Second, a calendar-time
cutoff absorbs every late-2022 shock: the BIS export controls
\cite{allen2022} fall eight weeks before ChatGPT and produce a statistically
indistinguishable estimate, so the design cannot name the shock. Third, the
panel is two aggregated units observed 46 times; with $\leq12$ quarters per
DiK cell, local slopes are estimated from few points on smooth curves, and
the full-sample 2021.5 kink shows that this setting manufactures nominal
$|z|>6$ where nothing happened.

The verdict is descriptive-only. The post-ChatGPT divergence in capex slopes
--- big-4 growth roughly doubling while the control aggregate's flattens ---
is a genuine feature of the data: sign-robust across every specification,
localized in 2023.0--2023.5 exactly as a lagged investment response to a
late-2022 shock would be, and visible per firm for three of the four
hyperscalers. But the design cannot certify it causally, and its headline
nominal $p\approx4\times10^{-6}$ overstates the evidence by roughly five
orders of magnitude relative to the placebo-calibrated $p$ of $0.125$--$0.25$.
Something bent hyperscaler capex around the turn of 2023; these data cannot
say it was ChatGPT.

\paragraph{Reproducibility.} All estimates were produced with \natex{} v0.2.0
\cite{natex}; estimation is deterministic (local weighted least squares), and
seed 0 was passed wherever the CLI accepts one. The EDGAR pulls are not
committed to the repository; the run of record retains the derived panel and
all estimate files. Figure~\ref{fig:fig1} regenerates deterministically from
the committed \texttt{figures/make\_fig.py}, which asserts the headline
estimates against the numbers of record before drawing.

\footnotesize
\begin{thebibliography}{9}
\setlength{\itemsep}{2pt}\setlength{\parskip}{0pt}

\bibitem{acemoglu2025}
Acemoglu, D. (2025).
The simple macroeconomics of AI.
\emph{Economic Policy}, 40(121), 13--58.

\bibitem{allen2022}
Allen, G.~C. (2022).
\emph{Choking Off China's Access to the Future of AI}.
Center for Strategic and International Studies, 11 October 2022.
\url{https://www.csis.org/analysis/choking-chinas-access-future-ai}

\bibitem{bertrand2004}
Bertrand, M., Duflo, E., and Mullainathan, S. (2004).
How much should we trust differences-in-differences estimates?
\emph{The Quarterly Journal of Economics}, 119(1), 249--275.

\bibitem{bockerman2025}
B\"ockerman, P., Jysm\"a, S., and Kanninen, O. (2025).
\emph{Difference-in-Kinks Design}.
IZA Discussion Paper No.~18313.
\url{https://docs.iza.org/dp18313.pdf}

\bibitem{card2015}
Card, D., Lee, D.~S., Pei, Z., and Weber, A. (2015).
Inference on causal effects in a generalized regression kink design.
\emph{Econometrica}, 83(6), 2453--2483.

\bibitem{eisfeldt2023}
Eisfeldt, A.~L., Schubert, G., and Zhang, M.~B. (2023).
Generative AI and firm values.
NBER Working Paper No.~31222.

\bibitem{ganong2018}
Ganong, P., and J\"ager, S. (2018).
A permutation test for the regression kink design.
\emph{Journal of the American Statistical Association}, 113(522), 494--504.

\bibitem{natex}
Hillebrandt, H. (2026).
\emph{natex: automated natural-experiment discovery and estimation} (version
0.2.0). Software.
\url{https://github.com/HaukeHillebrandt/natex}

\end{thebibliography}

\end{document}
